\subsection{発展的な読み物}
\par\smallskip\noindent\refstepcounter{table}\label{tab:ch17-23}表 \thetable: 発展的な読み物における文献・読むと得られること\par\smallskip
表 \thetable は、発展的な読み物における文献・読むと得られることを比較・確認しやすい形で整理したものである。\par\smallskip
\begin{longtable}{>{\raggedright\arraybackslash}p{0.26\linewidth} >{\raggedright\arraybackslash}p{0.68\linewidth}}
\toprule
文献 & 読むと得られること \\
\midrule
\endfirsthead
\toprule
文献 & 読むと得られること \\
\midrule
\endhead
Rosen, Discrete Mathematics and Its Applications & 論理、証明、集合、関係、関数、グラフ、木、数論、数え上げ、離散確率を体系的に学べる。ソフトウェアの基礎数学として最初に読む価値が高い。 \\
Cheney and Kincaid, Numerical Mathematics and Computing & 数値計算、近似、誤差、数値アルゴリズムを学べる。シミュレーション、最適化、科学技術計算に関わる場合に有用である。 \\
\bottomrule
\end{longtable}
